\section{What has been resolved, what remains, and how to reuse the work}
\label{sec:synthesis}

\subsection{A common framework with substantive domain theorems}

The unifying object is a \emph{specified presentation}, not a discrete
replacement of the mathematical object. A name may represent a point in an
infinite-dimensional Hilbert space, a function on an arbitrary presented
compact metric space, or a germ of a holomorphic function. What is finite
is a requested observation and the computation proving its error bound.
Four recurring ingredients determine the successful cases:
\begin{enumerate}
\item A dense family of finite objects with primitive-recursive operations.
\item A modulus converting the desired output precision to input demands.
\item A stability inequality converting approximate local work into a
coherent global result.
\item A semantic theorem identifying the represented result with the
intended continuous object.
\end{enumerate}
None of the four is automatic. The report produces them for concrete
classes rather than adding them as names for the desired conclusion.

\begin{longtable}{p{.25\textwidth}p{.36\textwidth}p{.30\textwidth}}
\toprule
Difficulty & Result supplied here & Exact remaining boundary\\
\midrule
\endfirsthead
\toprule
Difficulty & Result supplied here & Exact remaining boundary\\
\midrule
\endhead
An operation is given only on dense finite data &
Theorem~\ref{thm:cm-extension} constructs its extension and proves independence
of names from a joint modulus. &
The dense evaluator and its joint modulus must actually be supplied.\\[2mm]
A change of presentation appears to change the object &
Proposition~\ref{prop:cm-products-density} constructs inverse name maps from
a dense isometry and an approximate inverse selector. &
A mere assertion of density does not uniformly provide that selector.\\[2mm]
Continuous formulas contain alternating quantifiers &
Theorems~\ref{thm:cm-backforth} and \ref{thm:cm-values} construct finite
approximations to their values, retaining alternation. &
Numerical approximation does not decide exact zero truth or prove
logical compactness.\\[2mm]
A family of continuous functions needs a common finite representation &
Theorem~\ref{thm:quant-ascoli} constructs tent interpolants and an explicit
entropy bound. Proposition~\ref{prop:ascoli-converse} gives a quantitative
converse. &
A common modulus and effective domain nets are required; pointwise
continuity alone is insufficient.\\[2mm]
Local formulas must be assembled without dividing by an uncertified
denominator &
Theorem~\ref{thm:guarded-gluing} derives a denominator lower bound, an
evaluator, a Lipschitz constant, and the gluing error. &
Approximate continuous gluing need not preserve holomorphicity.\\[2mm]
Finite approximations need to define one solution &
Theorem~\ref{thm:hilbert-solver} compares two finite runs directly.
The minimax example also produces a coherence modulus from a proved gap. &
Unstable approximate solutions need not admit any effective coherent choice.\\[2mm]
An infinite-dimensional equation needs an actual solution producer &
Theorem~\ref{thm:hilbert-solver} derives the contraction gap, clock and
precision; Theorem~\ref{thm:pde} constructs a weak elliptic solution. &
The strong operator data and positive coercivity margin are essential.\\[2mm]
Weak Hilbert data lack a norm name &
Theorem~\ref{thm:tail-equivalence} identifies the exact tail information
needed; Proposition~\ref{prop:riesz-obstruction} proves it cannot be
manufactured in general. &
No blanket Riesz or weak-data Lax--Milgram algorithm follows.\\[2mm]
Differentiation consumes uncontrolled precision &
Theorems~\ref{thm:analytic-tail} and \ref{thm:analytic-entropy} construct
simultaneous derivative approximants with explicit degree and precision. &
A radius gap and coefficient bound are retained; there is no same-disk
uniform differentiation modulus.\\
\bottomrule
\end{longtable}

\subsection{The source audit changes some apparent problems}

The review found that several general proof-transport ideas are already
implemented in the current source, even where old packet README files still
describe uncompiled proposals. Conversely, not every declaration located
under \sourcepath{Derivations/} has an object-theory conclusion.
The inspected \sourcepath{Foundations/FiniteUseSignObstruction.lean}
expressly concerns external Lean propositions about oracle computation.
\sourcepath{Foundations/CoherenceObstruction.lean} supplies finite
arithmetic diagonal derivations while explicitly leaving its full
simulator-index noncomputability argument outside the formalized result.
The source ledger records these distinctions. Merely counting files,
theorem declarations, or compilation successes would obscure the real
verification boundary.

The work therefore gives priority to new mathematical instances and to
closing concrete representation obligations. A new general-purpose
proof-graph wrapper would have contributed less than the explicit PDE
operator construction, the continuous function-space cover, or the
tail-modulus equivalence.

\subsection{Formal and executable deliverables}

\begin{center}
\small
\begin{tabular}{p{.29\textwidth}p{.62\textwidth}}
\toprule
Artifact & Status and purpose\\
\midrule
\sourcepath{report.tex}, included sections, and \sourcepath{report.pdf} &
Full human mathematical proofs, logical interpretation, source discussion,
and formalization plan.\\
\sourcepath{lean/QuantitativeKernel.lean} &
19 fully proved host Lean lemmas, compiled with Lean 4.32.0; no
\texttt{sorry} and no newly declared axioms.\\
\sourcepath{lean/Audit.lean} and \sourcepath{lean/verification.log} &
Actual axiom audit and replay output. Some rational lemmas inherit standard
Lean metatheoretic axioms, including \texttt{Classical.choice}.\\
\sourcepath{lean/OBJECT_LANGUAGE_PLAN.md} &
Exact denominator-cleared invariants, source interfaces and remaining
\texttt{Prf PRAAx} targets.\\
\sourcepath{code/} &
Four standard-library Python programs using exact rational arithmetic.\\
\sourcepath{evidence/} &
Source audit, finite check results, and document/build verification.\\
\bottomrule
\end{tabular}
\end{center}

The finite examples include the alternating minimax formula and its optimizer;
continuous interpolation, clipping and guarded gluing; a genuinely
infinite diagonal Hilbert equation evaluated through finite support;
the variable-coefficient elliptic equation of Example~\ref{ex:pde}; and
complex derivative jets. Their success is evidence about implementations
and constants. The universal theorems rest on the printed mathematical
proofs, and their fully internal formalizations remain the enumerated targets.

\subsection{An ordered path into the existing repository}

The most useful first addition to the object library is the
denominator-cleared perturbed-contraction and clock package.
It has a small arithmetic dependency surface and immediately supports the
Hilbert, analytic-tail and stability applications.
Next come rational polynomial energy operations and the finite
two-run comparison. The operator and PDE endpoints can then be stated in
the same named-object style as existing analysis derivations.

In parallel, the partition-of-unity rational evaluator can be built on the
existing finite-net and rational-arithmetic infrastructure. That supplies
a reusable function-space layer. A continuous-formula module should
then implement fixed syntax and finite quantifier restrictions explicitly,
instead of treating a finite net as a finite submodel. Finally, the
split-radius complex estimates provide a specialized \(C^J\) family
presentation to which both function-space and continuous-logic observations
can be applied.

Each integration step should end with an explicit closed arithmetic
sentence or open contract and a kernel-checked
\(\texttt{Prf PRAAx}\) term, followed by an axiom audit.
The semantic completion or PDE explanation belongs beside that endpoint,
so a future user can see both what was proved internally and what continuous
mathematics it represents.
